% ============================================================================
% ARBEITSBLATT DS 8 - Gruppe A: Mi persona favorita
% ============================================================================

\documentclass[11pt,a4paper]{article}

\usepackage[utf8]{inputenc}
\usepackage[T1]{fontenc}
\usepackage[ngerman]{babel}
\usepackage[left=2cm,right=2cm,top=1.8cm,bottom=1.5cm]{geometry}
\usepackage{helvet}
\usepackage[table]{xcolor}
\usepackage{tabularx}
\usepackage{array}
\usepackage{enumitem}
\usepackage{tcolorbox}
\usepackage{fancyhdr}
\usepackage{lastpage}
\usepackage{amssymb}

\renewcommand{\familydefault}{\sfdefault}

\definecolor{spanisch}{HTML}{c41e3a}
\definecolor{grau}{HTML}{666666}
\definecolor{hellgrau}{HTML}{f5f5f5}

\pagestyle{fancy}
\fancyhf{}
\fancyhead[L]{\small\textcolor{grau}{Spanisch Spätbeginner -- Gruppe A}}
\fancyhead[R]{\small\textcolor{grau}{AB-08-A}}
\fancyfoot[C]{\small\thepage/\pageref{LastPage}}
\renewcommand{\headrulewidth}{0.5pt}

\newcommand{\luecke}[1]{\underline{\hspace{#1}}}
\newcommand{\punktebox}[1]{\hfill\fbox{\small \_\_/#1}}
\newcommand{\es}[1]{\textit{#1}}

\newtcolorbox{merkkasten}[1][]{
    colback=white,
    colframe=spanisch,
    fonttitle=\bfseries\color{spanisch},
    title=#1,
    boxrule=1.5pt,
    arc=0pt,
    left=5pt, right=5pt, top=5pt, bottom=5pt
}

\newtcolorbox{tippbox}{
    colback=hellgrau,
    colframe=grau,
    boxrule=0.5pt,
    arc=0pt,
    left=5pt, right=5pt, top=3pt, bottom=3pt
}

\newcounter{aufgabe}
\newcommand{\aufgabe}[2]{%
    \stepcounter{aufgabe}%
    \vspace{0.8em}%
    \noindent\textbf{\theaufgabe.} #1 \punktebox{#2}%
    \vspace{0.3em}%
}

\begin{document}

% ============================================================================
% KOPFBEREICH
% ============================================================================
\noindent
\begin{tabularx}{\textwidth}{@{}Xr@{}}
    {\Large\textbf{\textcolor{spanisch}{Mi persona favorita}}} & Name: \luecke{5cm} \\[0.3em]
    {\small DS 8 -- Creatividad y Juego} & Datum: \luecke{3cm} \\
\end{tabularx}

\vspace{0.3em}
\hrule
\vspace{0.8em}

% ============================================================================
% MERKKASTEN: Adjektive zur Personenbeschreibung
% ============================================================================
\begin{merkkasten}[Kasten 1: Adjetivos para describir personas]
\vspace{0.2em}

\begin{tabularx}{\textwidth}{@{}|X|X|X|@{}}
\hline
\rowcolor{hellgrau}
\textbf{Positiv} & \textbf{Neutral} & \textbf{Charakter} \\
\hline
simpático/-a & tranquilo/-a & trabajador/-a \\
(sympathisch) & (ruhig) & (fleißig) \\[0.3em]
\hline
inteligente & serio/-a & responsable \\
(intelligent) & (ernst) & (verantwortungsvoll) \\[0.3em]
\hline
divertido/-a & reservado/-a & organizado/-a \\
(lustig) & (zurückhaltend) & (organisiert) \\[0.3em]
\hline
amable & \luecke{2.5cm} & creativo/-a \\
(freundlich) & \luecke{2.5cm} & (kreativ) \\[0.3em]
\hline
paciente & \luecke{2.5cm} & \luecke{2.5cm} \\
(geduldig) & \luecke{2.5cm} & \luecke{2.5cm} \\[0.3em]
\hline
\end{tabularx}

\vspace{0.3em}
\end{merkkasten}

\vspace{0.8em}

% ============================================================================
% MERKKASTEN 2: Strukturen
% ============================================================================
\begin{merkkasten}[Kasten 2: Satzstrukturen für Personenbeschreibungen]
\vspace{0.2em}

\begin{tabularx}{\textwidth}{@{}|p{4cm}|X|@{}}
\hline
\rowcolor{hellgrau}
\textbf{Struktur} & \textbf{Beispiel} \\
\hline
\textbf{ser + Adjektiv} & Mi madre \textbf{es} muy amable. \\[0.3em]
\hline
\textbf{tener + años} & \textbf{Tiene} 45 años. \\[0.3em]
\hline
\textbf{Possessivpronomen} & \textbf{Mi} hermano, \textbf{su} amiga, \textbf{nuestro} profesor \\[0.3em]
\hline
\textbf{Konnektoren} & también, además, pero, porque \\[0.3em]
\hline
\end{tabularx}

\vspace{0.3em}
\end{merkkasten}

\vspace{1em}

% ============================================================================
% AUFGABEN
% ============================================================================

\aufgabe{Wortschatz-Staffel: Schreibe so viele Adjektive wie möglich. (2 Min.)}{6}

\begin{tippbox}
\textbf{Ziel:} Mindestens 10 verschiedene Adjektive zur Personenbeschreibung.
\end{tippbox}

\vspace{0.3em}
\begin{tabularx}{\textwidth}{@{}|X|X|X|X|X|@{}}
\hline
1. \luecke{2cm} & 2. \luecke{2cm} & 3. \luecke{2cm} & 4. \luecke{2cm} & 5. \luecke{2cm} \\[0.6em]
\hline
6. \luecke{2cm} & 7. \luecke{2cm} & 8. \luecke{2cm} & 9. \luecke{2cm} & 10. \luecke{2cm} \\[0.6em]
\hline
\end{tabularx}

\vspace{0.8em}

\aufgabe{Ergänze \es{ser} in der richtigen Form.}{6}

\begin{enumerate}[label=\alph*), leftmargin=2em]
    \item Mi padre \luecke{2cm} muy trabajador.
    \item Mis abuelos \luecke{2cm} simpáticos y pacientes.
    \item ¿Cómo \luecke{2cm} tu mejor amigo?
    \item Nosotros \luecke{2cm} estudiantes creativos.
    \item Ella \luecke{2cm} inteligente y responsable.
    \item Yo \luecke{2cm} una persona organizada.
\end{enumerate}

\vspace{0.8em}

\aufgabe{Verbinde mit Konnektoren. Schreibe vollständige Sätze.}{8}

\begin{tippbox}
\textbf{Verwende:} \es{también} (auch), \es{además} (außerdem), \es{pero} (aber), \es{porque} (weil)
\end{tippbox}

\vspace{0.3em}
\begin{enumerate}[label=\alph*), leftmargin=2em]
    \item Mi hermana es simpática. / Es muy inteligente.

    \vspace{0.3cm}
    \luecke{13cm}

    \item Mi profesor es serio. / Es muy amable.

    \vspace{0.3cm}
    \luecke{13cm}

    \item Me gusta mi abuela. / Es muy paciente.

    \vspace{0.3cm}
    \luecke{13cm}

    \item Mi amigo es divertido. / Es creativo.

    \vspace{0.3cm}
    \luecke{13cm}
\end{enumerate}

\vspace{0.8em}

\aufgabe{Poster-Vorbereitung: Beschreibe deine Lieblingsperson.}{8}

\begin{tippbox}
\textbf{Schreibe mindestens 5 Sätze} für dein Poster. Nutze verschiedene Strukturen!
\end{tippbox}

\vspace{0.5em}
\textbf{Struktur für dein Poster:}

\vspace{0.3em}
\begin{tabularx}{\textwidth}{@{}|X|@{}}
\hline
\textbf{1. Wer ist die Person?} \\[0.3em]
Mi persona favorita es \luecke{10cm} \\[0.5em]
Es mi \luecke{4cm} (madre, padre, amigo/-a, profesor/-a...) \\[0.6em]
\hline
\textbf{2. Wie ist sie? (mind. 3 Adjektive)} \\[0.3em]
\luecke{14cm} \\[0.6em]
\luecke{14cm} \\[0.6em]
\hline
\textbf{3. Warum magst du sie?} \\[0.3em]
Me gusta mucho porque \luecke{9cm} \\[0.6em]
\luecke{14cm} \\[0.6em]
\hline
\end{tabularx}

\vspace{0.8em}

\aufgabe{Gallery Walk: Bewerte 3 Poster deiner Mitschüler.}{4}

\begin{tabularx}{\textwidth}{@{}|p{4cm}|c|c|c|p{4cm}|@{}}
\hline
\rowcolor{hellgrau}
\textbf{Name / Poster} & \textbf{Sprache} & \textbf{Kreativität} & \textbf{Gesamt} & \textbf{Notiz} \\
\hline
& $\square\square\square$ & $\square\square\square$ & $\square\square\square$ & \\[0.8em]
\hline
& $\square\square\square$ & $\square\square\square$ & $\square\square\square$ & \\[0.8em]
\hline
& $\square\square\square$ & $\square\square\square$ & $\square\square\square$ & \\[0.8em]
\hline
\end{tabularx}

\vspace{0.3em}
{\small Legende: $\square$ = 1 Punkt (max. 3 pro Kategorie)}

\vspace{1em}
\hrule
\vspace{0.3em}
\begin{flushright}
\textbf{Gesamt: \luecke{1cm} / 32 Punkte}
\end{flushright}

\end{document}
